\begin{abstractCH}

近年來，非地球同步軌道（non-geostationary satellite orbit, NGSO）衛星系統，特別是低地球軌道（low Earth orbit, LEO）衛星星座，已成為未來寬頻衛星通訊與非地面網路的重要發展方向。然而，當 LEO 衛星下行傳輸與既有地球同步軌道（geostationary satellite orbit, GSO）系統共用相同或相鄰頻段時，可能對 GSO 地面站造成同頻干擾。為保護 GSO 系統，NGSO 系統必須遵守等效功率通量密度（equivalent power flux-density, EPFD）限制。當 LEO 衛星通過 GSO 地面站附近，且與 GSO 方向形成近似共線幾何時，部分高干擾風險波束可能需要被關閉以滿足 EPFD 限制。然而，波束關閉雖然能有效降低干擾，卻會使原本由關閉波束服務的使用者失去服務，進而影響 LEO 使用者的服務連續性與滿足率。

為解決上述問題，本文提出 Safe-Beam Assisted Power-Release Relay（SAPR-R），一種適用於 NGSO--GSO 共存情境下多波束 LEO 衛星系統的分階段中繼式服務恢復方法。SAPR-R 利用 LEO 衛星軌道與波束覆蓋範圍可預測的特性，預先建立 critical satellite、helper satellites、closed beams、safe beams，以及 critical/helper 衛星之間的波束重疊對應關係。當 critical satellite 為滿足 EPFD 限制而關閉部分高風險波束時，SAPR-R 先透過 critical satellite 上仍可運作的 safe beams 接手部分 helper users，以釋放 helper satellite 的可用功率；接著將釋放出的功率分配給 helper relay beams，以提升其接手 closed-beam users 的能力；最後，透過 beam-pair local relay 將受 critical-beam shutdown 影響的使用者重新分配至可行的 helper relay beams，完成服務恢復。

本文透過 STK 與 MATLAB 建立模擬環境，評估 SAPR-R 在 critical satellite 通過 GSO 地面站附近期間的 EPFD 合規性與使用者服務表現。模擬結果顯示，SAPR-R 能使 aggregate EPFD 維持在規範限制以下，並相較於 baseline methods 達到更高的平均使用者滿足率。此外，SAPR-R 能提升成功被 helper relay beams 接手的 closed-beam users 數量，並降低低滿足率使用者的比例。整體而言，SAPR-R 能在滿足 EPFD 限制的同時，有效改善 LEO 使用者的服務連續性，並避免在每個 time slot 中求解高複雜度的全域最佳化問題，提供一種兼具實用性與可行性的中繼式服務恢復方法。

關鍵詞：NGSO--GSO 共存、EPFD、LEO 衛星、多波束衛星系統、服務恢復、中繼、非地面網路、SAPR-R

\end{abstractCH}
